\begin{example}
\label{editorial-source-example-product-matrices-zero}
Let $k$ be a field. The space $\Spec(k[x, y]/(xy))$
has two irreducible components: namely the $x$-axis and the
$y$-axis. As a generalization, let
$$
R = k[x_{11}, x_{12}, x_{21}, x_{22}, y_{11}, y_{12}, y_{21}, y_{22}]/
\mathfrak a,
$$
where $\mathfrak a$ is the ideal in
$k[x_{11}, x_{12}, x_{21}, x_{22}, y_{11}, y_{12}, y_{21}, y_{22}]$
generated by the entries of the $2 \times 2$ product matrix
$$
\left(
\begin{matrix}
x_{11} & x_{12}\\
x_{21} & x_{22}
\end{matrix}
\right)
 \left(
\begin{matrix}
y_{11} & y_{12}\\
y_{21} & y_{22}
\end{matrix}
\right).
$$
In this example we will describe $\Spec(R)$.

\medskip\noindent
To prove the statement about $\Spec(k[x, y]/(xy))$ we argue as follows.
If $\mathfrak p\subset k[x,y]$ is any prime ideal containing $xy$,
then either $x$ or $y$ belongs to $\mathfrak p$. Hence the minimal such
prime ideals are just $(x)$ and $(y)$. In case $k$ is
algebraically closed, the $\text{max-Spec}$ of these components
can then be visualized as the point sets of $y$- and $x$-axis.

\medskip\noindent
For the generalization, note that we may identify the closed
points of the spectrum of
$k[x_{11}, x_{12}, x_{21}, x_{22}, y_{11}, y_{12}, y_{21}, y_{22}]$
with the space of matrices
$$
\left\{ (X, Y) \in \text{Mat}(2, k)\times \text{Mat}(2, k) \mid
X = \left(
\begin{matrix}
x_{11} & x_{12}\\
x_{21} & x_{22}
\end{matrix}
\right),
Y= \left(
\begin{matrix}
y_{11} & y_{12}\\
y_{21} & y_{22}
\end{matrix}
\right)
\right\}
$$
if $k$ is algebraically closed. We assume this for the following
description in terms of closed $k$-points. The coordinate-ring proof
below establishes the component statement over every field $k$.
Now define a group action of
$\text{GL}(2, k)\times \text{GL}(2, k)\times \text{GL}(2, k)$
on the space of matrices $\{(X, Y)\}$ by
$$
(g_1, g_2, g_3) \times (X, Y) \mapsto ((g_1Xg_2^{-1}, g_2Yg_3^{-1})).
$$
Here, also observe that the algebraic set
$$
\text{GL}(2, k)\times \text{GL}(2, k)\times \text{GL}(2, k) \subset
\text{Mat}(2, k)\times \text{Mat}(2, k) \times \text{Mat}(2, k)
$$
is irreducible since it is the max spectrum of the domain
$$
k[x_{11}, x_{12}, \ldots, z_{21}, z_{22}, (x_{11}x_{22}-x_{12}x_{21})^{-1}
, (y_{11}y_{22}-y_{12}y_{21})^{-1}, (z_{11}z_{22}-z_{12}z_{21})^{-1}].
$$
Since the image of an irreducible algebraic set is still
irreducible, the orbit closures below are irreducible. There are
six orbits, indexed by the pairs of ranks $(r,s)$ with $r+s\leq2$:
$$
(2,0),\quad(1,1),\quad(0,2),\quad(1,0),\quad(0,1),\quad(0,0).
$$
Indeed $XY=0$ says $\operatorname{Im}(Y)\subset\ker(X)$.
Choose a basis of the middle space starting with a basis of
$\operatorname{Im}(Y)$, extend it to $\ker(X)$ and then to the
whole middle space. Lift a basis of $\operatorname{Im}(Y)$ through
$Y$ and extend those lifts by a basis of $\ker(Y)$ in the source.
In the target extend the images under $X$ of the remaining middle
basis vectors to a basis. These choices show that the ranks determine
the orbit. If their basis matrices in the original coordinates are
$C_1,C_2,C_3$, respectively, then
$g_1=C_1^{-1}$, $g_2=C_2^{-1}$, $g_3=C_3^{-1}$ give exactly
the displayed action, with transformed matrices
$C_1^{-1}XC_2$ and $C_2^{-1}YC_3$. It preserves the original
equations because the transformed product is $g_1XYg_3^{-1}$,
and the inverse transformation uses $(g_1^{-1},g_2^{-1},g_3^{-1})$.
The three cases according to the rank of $X$ are:
\begin{enumerate}
\item $\exists (g_1, g_2)$ such that $g_1Xg_2^{-1} = I_{2\times 2}$.
Then $Y$ is necessarily $0$, which as an algebraic set is
invariant under the group action. It follows that this orbit is
contained in the irreducible algebraic set defined by the prime
ideal $(y_{11}, y_{12}, y_{21}, y_{22})$. Taking the closure, we see
that $(y_{11}, y_{12}, y_{21}, y_{22})$ is actually a component.
\item $\exists (g_1, g_2)$ such that
$$
g_1Xg_2^{-1} = \left(
\begin{matrix}
1 & 0 \\
0 & 0
\end{matrix}
\right).
$$
This case occurs if and only if $X$ is a rank 1 matrix,
and furthermore, $Y$ is killed by such an $X$ if and only if
$$
x_{11}y_{11}+x_{12}y_{21} = 0; \quad x_{11}y_{12}+x_{12}y_{22} = 0;
$$
$$
x_{21}y_{11}+x_{22}y_{21} = 0; \quad x_{21}y_{12}+x_{22}y_{22} = 0.
$$
Fix a rank $1$ matrix $X$. The nonzero matrices $Y$ satisfying these
equations form an irreducible locally closed set, for the following
reason (the pair $(X,0)$ belongs to the closure of the first orbit):
$0 = g_1Xg_2^{-1}g_2Y$ implies that
$$
g_2Y = \left(
\begin{matrix}
0 & 0 \\
y_{21}' & y_{22}'
\end{matrix}
\right).
$$
With a further $\text{GL}(2, k)$-action on the right by $g_3$,
$g_2Y$ can be brought into
$$
g_2Yg_3^{-1} = \left(
\begin{matrix}
0 & 0 \\
0 & 1
\end{matrix}
\right),
$$
and thus such $Y$'s form an irreducible algebraic set
isomorphic to the image of $\text{GL}(2, k)$ under this action. Finally,
notice that the ``rank 1" condition for $X$'s forms an open dense
subset of the irreducible algebraic set
$\det X = x_{11}x_{22} - x_{12}x_{21} = 0$.
On the open set where both matrices are nonzero, the four displayed
product equations and $\det X=0$ describe the rank-$(1,1)$ orbit.
To identify its closure, retain the polynomial ring
$B=k[x_{11},x_{12},x_{21},x_{22},y_{11},y_{12},y_{21},y_{22}]$
and put
$$
\begin{aligned}
\mathfrak q_M=(
&x_{11}y_{11}+x_{12}y_{21},\ x_{11}y_{12}+x_{12}y_{22},\\
&x_{21}y_{11}+x_{22}y_{21},\ x_{21}y_{12}+x_{22}y_{22},\\
&x_{11}x_{22}-x_{12}x_{21},\ y_{11}y_{22}-y_{12}y_{21}).
\end{aligned}
$$
The proof below shows that this is prime and cuts out that closure.
The sixth equation is already a consequence of the first four on
the open set where $X$ is nonzero: writing $a_{ij}=(XY)_{ij}$,
the exact matrix identity $X\det(Y)=(XY)\operatorname{adj}(Y)$ gives
$$
\begin{aligned}
x_{11}\det Y&=a_{11}y_{22}-a_{12}y_{21},&
x_{12}\det Y&=-a_{11}y_{12}+a_{12}y_{11},\\
x_{21}\det Y&=a_{21}y_{22}-a_{22}y_{21},&
x_{22}\det Y&=-a_{21}y_{12}+a_{22}y_{11}.
\end{aligned}
$$
Thus on each standard open where an entry of $X$ is invertible, the
original five-equation ideal and $\mathfrak q_M$ give the same quotient.
No sixth equation is discarded in the closure.
\item $\exists (g_1, g_2)$ such that $g_1Xg_2^{-1} = 0$, or
equivalently, $X = 0$. Then $Y$ can be arbitrary and this component
is thus defined by $(x_{11}, x_{12}, x_{21}, x_{22})$.
\end{enumerate}

\medskip\noindent
Here is the coordinate-ring proof, valid over the original arbitrary
field $k$. Keep $B$ and the four-entry ideal $\mathfrak a$ of the
original product $XY$, so $R=B/\mathfrak a$. Put
$$
\mathfrak q_X=(x_{11},x_{12},x_{21},x_{22}),\qquad
\mathfrak q_Y=(y_{11},y_{12},y_{21},y_{22}),
$$
and retain $\mathfrak q_M=\mathfrak a+(\det X,\det Y)$ as displayed
above. The first two ideals are prime, since their quotient rings
are polynomial rings in the four remaining entries.

For the middle ideal use the exact rearrangement of the original entries
$$
M=\begin{pmatrix}
x_{11}&x_{21}&y_{21}&y_{22}\\
x_{12}&x_{22}&-y_{11}&-y_{12}
\end{pmatrix}.
$$
Its six $2\times2$ minors, in column order
$(12),(13),(14),(23),(24),(34)$, are
$$
\det X,\quad-a_{11},\quad-a_{12},\quad-a_{21},\quad-a_{22},\quad\det Y.
$$
The entries of $M$ are eight independent polynomial coordinates:
the inverse rearrangement recovers each original entry, including
$y_{11}=-M_{23}$ and $y_{12}=-M_{24}$. Define
$$
\Phi:B\longrightarrow k[\lambda,\mu,c_1,c_2,c_3,c_4]
$$
by
$$
\begin{aligned}
x_{11}&\mapsto\lambda c_1,&x_{12}&\mapsto\mu c_1,&
x_{21}&\mapsto\lambda c_2,&x_{22}&\mapsto\mu c_2,\\
y_{11}&\mapsto-\mu c_3,&y_{12}&\mapsto-\mu c_4,&
y_{21}&\mapsto\lambda c_3,&y_{22}&\mapsto\lambda c_4.
\end{aligned}
$$
Thus the two rows of $M$ map to $(\lambda c_i)_i$ and
$(\mu c_i)_i$. Every minor maps to zero, proving
$\mathfrak q_M\subset\ker\Phi$.

To prove equality, write $s_i=M_{1i}$ and $t_i=M_{2i}$ for
$1\leq i\leq4$. The image of a monomial
$\prod_i s_i^{u_i}t_i^{v_i}$ is
$\lambda^{\sum_i u_i}\mu^{\sum_i v_i}\prod_i c_i^{u_i+v_i}$.
Two monomials have the same image exactly when all their column
totals $u_i+v_i$ agree and their first-row totals $\sum_i u_i$ agree.
For two such monomials with first-row exponents $u_i,u'_i$, if
$u\ne u'$ choose $i,j$ with $u_i>u'_i$ and $u_j<u'_j$.
Then $s_i$ and $t_j$ divide the first monomial. Replacing the factor
$s_it_j$ by $t_is_j$ changes it by a multiple of the minor
$s_it_j-t_is_j$, preserves all row and column totals, and decreases
$\sum_i|u_i-u'_i|$ by $2$. Repetition reaches the second monomial,
proving their difference belongs to the minor ideal.
For an arbitrary polynomial in $\ker\Phi$, group its finitely many
terms according to their image monomial. Distinct image monomials
are linearly independent over $k$, so the sum of coefficients in
each group is zero. Choosing one monomial in each group expresses
that whole group as a $k$-linear combination of the differences
just proved to lie in the minor ideal. Consequently
$\ker\Phi=\mathfrak q_M$. The induced map
$B/\mathfrak q_M\hookrightarrow k[\lambda,\mu,c_1,c_2,c_3,c_4]$
is injective, so $\mathfrak q_M$ is prime over every field,
with all original signs retained also in characteristic $2$.

Now let $\mathfrak p\subset B$ be any prime containing
$\mathfrak a$. In the field $\operatorname{Frac}(B/\mathfrak p)$,
the original matrices satisfy $XY=0$. If $\det X\notin\mathfrak p$,
then $X$ is invertible over this field and $Y=0$, so
$\mathfrak q_Y\subset\mathfrak p$. If
$\det Y\notin\mathfrak p$, the same argument gives
$\mathfrak q_X\subset\mathfrak p$. If both determinants belong
to $\mathfrak p$, then $\mathfrak q_M\subset\mathfrak p$.
These three prime ideals are pairwise incomparable.
For example $x_{11}\in\mathfrak q_X\setminus\mathfrak q_M$
since $\Phi(x_{11})=\lambda c_1\ne0$, while
$\det Y\in\mathfrak q_M\setminus\mathfrak q_X$ since it remains
a nonzero polynomial in the quotient by $\mathfrak q_X$.
The analogous facts with $X,Y$ exchanged prove incomparability
with $\mathfrak q_Y$, and the two coordinate ideals are
incomparable by their distinct sets of variables.
Therefore these are exactly the three minimal primes over
$\mathfrak a$. Their images in $R$ define exactly the three
irreducible components of $\Spec(R)$.

Finally, over an algebraically closed field the rank-$(2,0)$ orbit
is the nonempty open $\det X\ne0$ in $V(\mathfrak q_Y)$ and is
dense there; the rank-$(0,2)$ orbit is dense in $V(\mathfrak q_X)$
by the identical argument. In the irreducible $V(\mathfrak q_M)$,
the open where both $X$ and $Y$ are nonzero is precisely the
rank-$(1,1)$ orbit. It is nonempty, as witnessed by
$$
X=\begin{pmatrix}1&0\\0&0\end{pmatrix},\qquad
Y=\begin{pmatrix}0&0\\0&1\end{pmatrix},
$$
and hence dense. Density of closed points in these affine finite
type spectra follows from the Jacobson results above, so these
are the same closures in the stated description by closed points.
The other three rank orbits lie in these components and give no
additional irreducible component.
\end{example}